% Source printed pp. 113-117. Subsections 4.1.7-4.1.8.
% Keep the indivisible table without stretching the preceding page's paragraphs.
\raggedbottom
\subsection{万有族}\label{src:4.1.7-subsection}
\begin{lemma}[广义 $2$-Yoneda 引理]\label{src:4.1.27}
设 $\mathcal X$ 是 $\Sch_{\text{ét}}$ 上的叠. 若 $\mathcal T$ 是代数叠, $U\to\mathcal T$ 是光滑呈示, 将范畴 $\mathcal X(\mathcal T)$ 定义为等化子
\[
\mathcal X(\mathcal T):=\operatorname{Eq}\left(
\begin{tikzcd}[column sep=small]
\mathcal X(U)\arrow[r,shift left=2pt]\arrow[r,shift right=2pt]&
\mathcal X(U\times_{\mathcal T}U)\arrow[r,shift left=3pt]\arrow[r]\arrow[r,shift right=3pt]&
\mathcal X(U\times_{\mathcal T}U\times_{\mathcal T}U)
\end{tikzcd}
\right),
\]
即对象是二元组 $(a,\alpha)$ 的数据, 其中 $a\in\mathcal X(U)$, $\alpha:p_1^*a\xrightarrow{\sim}p_2^*a$ 是满足上循环条件 $p_{23}^*\alpha\circ p_{12}^*\alpha=p_{13}^*\alpha$ 的同构, 而态射 $(a,\alpha)\to(a',\alpha')$ 是满足 $p_2^*\beta\circ\alpha=\alpha'\circ p_1^*\beta$ 的态射 $\beta:a\to a'$ 的数据.
存在自然的范畴等价
\begin{equation}\label{src:4.1.28}
\Mor(\mathcal T,\mathcal X)\longrightarrow\mathcal X(\mathcal T).
\end{equation}
特别地, $\mathcal X(\mathcal T)$ 与呈示的选择无关.
\end{lemma}
\begin{proof}
记光滑呈示为 $q:U\to\mathcal T$, 考虑 $2$-交换图
\[
\begin{tikzcd}
U\times_{\mathcal T}U\arrow[r,"p_2"]\arrow[d,"p_1"']&U\arrow[d,"q"]\\
U\arrow[r,"q"']&\mathcal T.
\arrow[Rightarrow,from=1-2,to=2-1,shorten <=1.6em,shorten >=1.6em,"\gamma"']
\end{tikzcd}
\]
给定态射 $f:\mathcal T\to\mathcal X$, 设 $a\in\mathcal X(U)$ 是通过 $2$-Yoneda 引理 (\bookref{3.4.21}) 对应于复合 $f\circ q:U\to\mathcal T\to\mathcal X$ 的对象.
将 $2$-同构 $\gamma:q\circ p_2\to q\circ p_1$ 与 $f$ 复合, 诱导 $2$-同构 $\alpha:p_1^*a\xrightarrow{\sim}p_2^*a$.
映射~\bookeqref{4.1.28}~将态射 $f$ 送到二元组 $(a,\alpha)\in\mathcal X(\mathcal T)$.

为证明~\bookeqref{4.1.28}~是等价, 先假设它对代数空间成立, 从而对 $U$, $U\times_{\mathcal T}U$ 和 $U\times_{\mathcal T}U\times_{\mathcal T}U$ 都成立.
为构造~\bookeqref{4.1.28}~的逆 $\Psi:\mathcal X(\mathcal T)\to\Mor(\mathcal T,\mathcal X)$, 设 $(a,\alpha)\in\mathcal X(\mathcal T)$, $f_a:U\to\mathcal X$ 是对应于 $a$ 的映射.
为定义叠态射 $\Psi(a,\alpha):\mathcal T\to\mathcal X$, 设 $b\in\mathcal T$ 是位于概形 $S$ 上的对象, 由映射 $f_b:S\to\mathcal T$ 分类, 考虑交换图
\[
\begin{tikzcd}[column sep=small]
S_2\arrow[r,hook]\arrow[drr]&S_1\arrow[r]&S_0\arrow[r]\arrow[d]\arrow[dr,phantom,"\square"]&U\arrow[r,"f_a"]\arrow[d,"q"]&\mathcal X\\
{}&{}&S\arrow[r,"f_b"']&\mathcal T,&{}
\end{tikzcd}
\]
其中 $S_0=S\times_{\mathcal T}U$ 是代数空间, $S_1\to S_0$ 是平展呈示, $S_2\to S$ 是通过 $S_1$ 分解的概形满射平展态射. 这样的态射存在, 因为光滑态射在平展局部具有截面 (推论~\bookref{A.3.5}).
顶行的复合定义 $\mathcal X$ 中位于 $S_2$ 上的对象 $c_2$.
$U\times_{\mathcal T}U$ 上的同构 $\alpha:p_1^*a\xrightarrow{\sim}p_2^*a$ 诱导 $S_2\times_S S_2$ 上的同构 $\beta:p_1^*c_2\xrightarrow{\sim}p_2^*c_2$.
由于 $\alpha$ 满足上循环条件, 且~\bookeqref{4.1.28}~对 $U\times_{\mathcal T}U\times_{\mathcal T}U$ 是等价, $\beta$ 也满足上循环条件. 因而存在 $\mathcal X$ 中位于 $S$ 上的对象 $c$, 拉回为 $c_2$.
令 $\Psi(a,\alpha)(b)=c\in\mathcal X(S)$. 验证这确实给出逆映射的细节留给读者.
最后, 当 $\mathcal T$ 是代数空间时, 使用平展呈示 $U\to\mathcal T$, 同样的论证说明~\bookeqref{4.1.28}~是等价.
\end{proof}
\begin{definition}[万有族]\label{src:4.1.29}
若 $\mathcal X$ 是代数叠, 则 $\mathcal X$ 上的万有族是等价~\bookeqref{4.1.28}~下对应于恒等态射 $\id_{\mathcal X}:\mathcal X\to\mathcal X$ 的对象 $u\in\mathcal X(\mathcal X)$ (在唯一同构的意义下唯一).
\end{definition}
\begin{exercise}[细节]\label{src:4.1.30}
设 $\mathcal X$ 是代数叠, $u\in\mathcal X(\mathcal X)$ 是万有族.
若 $g:\mathcal S\to\mathcal T$ 是代数叠态射, 证明存在自然的拉回函子 $g^*:\mathcal X(\mathcal T)\to\mathcal X(\mathcal S)$; 并证明若从概形出发的映射 $f_a:T\to\mathcal X$ 由对象 $a\in\mathcal X(T)$ 分类, 则 $a$ 同构于拉回 $f_a^*u$.
\end{exercise}
实际使用中, 对来自模问题的代数叠 $\mathcal X$, $\mathcal X(\mathcal S)$ 中对象的几何意义通常很清楚.
\begin{example}[$\mathcal M_g$ 的万有族]\label{src:4.1.31}
代数叠 $\mathcal T$ 上的光滑曲线族 $\mathcal C\to\mathcal T$ 是可由概形表出, 固有, 平坦且有限表现的代数叠态射, 使得对每个几何点 $\Spec k\to\mathcal T$, 纤维 $C_s:=\mathcal C\times_{\mathcal S}k$ 是亏格 $g$ 的光滑, 连通的射影曲线.
下降理论将 $\mathcal M_g(\mathcal T)$ 与 $\mathcal T$ 上的光滑曲线族组成的范畴等同.
设 $\mathcal U_g\to\mathcal M_g$ 为对应于万有族 $u\in\mathcal M_g(\mathcal M_g)$ 的曲线族.
根据习题~\bookref{3.4.42}, 万有族可与忘掉截面的映射 $\mathcal M_{g,1}\to\mathcal M_g$ 等同, 其中 $\mathcal M_{g,1}$ 是光滑单标点曲线的叠.
对对应于光滑曲线族 $C\to S$ 的每个态射 $S\to\mathcal M_g$, 有 Cartesian 图
\[
\begin{tikzcd}
C\arrow[r]\arrow[d]\arrow[dr,phantom,"\square"]&\mathcal U_g=\mathcal M_{g,1}\arrow[d]\\
S\arrow[r]&\mathcal M_g.
\end{tikzcd}
\]
关于 $\mathcal U_g\to\mathcal M_g$ 的直观图像, 见图~\bookref{1.3.27}.
\end{example}
\begin{example}[$\cBun_{r,d}(C)$ 的万有族]\label{src:4.1.32}
后面将利用拟凝聚层理论定义代数叠上的向量丛 (见 \S\bookref{7.1}).
由下降理论, $\cBun_{r,d}(C)(\mathcal T)$ 与 $C\times\mathcal T$ 上秩为 $r$, 次数为 $d$ 的向量丛的群胚等同.
万有族是 $C\times\cBun_{r,d}(C)$ 上的向量丛 $\mathcal E_{\mathrm{univ}}$, 具有如下性质: 若 $S$ 是概形, $E$ 是 $C\times S$ 上由映射 $f:S\to\cBun_{r,d}(C)$ 分类的向量丛, 则 $E\cong(\id\times f)^*\mathcal E_{\mathrm{univ}}$.
\end{example}

\subsection{目标与展望}\label{src:4.1.8-subsection}
后续各章将发展代数空间与代数叠的基础理论, 但我们先介绍其中的一些要点.

\textbf{对角线的重要性.} 听到别人讨论代数叠时, 你也许疑惑过, 为什么对角线会引起如此多的关注. 原因就在于: 对角线编码了叠的特性!

首先, 代数叠的对角线 $\mathcal X\to\mathcal X\times\mathcal X$ 可表, 而代数空间的对角线 $X\to X\times X$ 可由概形表出 (定理~\bookref{4.2.1}).
域值点 $x:\Spec k\to\mathcal X$ 的自同构群或稳定子定义为层
\[
G_x:=\underline{\Aut}_{\mathcal X(k)}(x)=\underline{\Isom}_{\mathcal X(k)}(x,x),
\]
根据习题~\bookref{3.4.39}, 它与纤维积 $\mathcal X\times_{\mathcal X\times\mathcal X}\Spec k$ 等同.
由对角线的可表性 (\bookref{4.2.1}), 稳定子 $G_x$ 可由 $k$ 上的群代数空间表出. 关于稳定子的进一步讨论, 见 \S\bookref{4.2.2-subsection}.

概形 (分别为分离概形) 的对角线是浸入 (分别为闭浸入).
代数空间的对角线未必是浸入, 而代数叠的对角线甚至未必是单态射, 因为在 $(x,x):\Spec k\to\mathcal X\times\mathcal X$ 上的纤维 (即稳定子 $G_x$) 可能非平凡.
定理~\bookref{4.6.4}~和~\bookref{4.6.6}~用对角线刻画代数空间与 Deligne--Mumford 叠, 如下表所示:
\begin{table}[H]
\centering
\begin{tabularx}{\textwidth}{*{3}{X}}
\toprule
空间类型&对角线的性质&稳定子的性质\\
\midrule
代数空间&单态射&平凡\\
Deligne--Mumford 叠&非分歧&约化的有限群\footnotemark\\
代数叠&任意&任意\\
\bottomrule
\end{tabularx}
\caption{代数空间与 Deligne--Mumford 叠的刻画. 关于模理论三分法的进一步图示说明, 见表~\bookref{1.0.1}, 图~\bookref{1.0.2}~及表~\bookref{1.6.5}.}\label{src:4.1.33}
\end{table}
\footnotetext{若对角线不拟紧, 则稳定子只能保证离散且约化.}
这一刻画将使我们能够推出 $\mathcal M_g$ 是 Deligne--Mumford 叠 (推论~\bookref{4.6.10}).

后面将在推论~\bookref{4.6.9}~和~\bookref{5.5.11}~中证明, 代数叠态射 $f:\mathcal X\to\mathcal Y$ 可表当且仅当对每个域值点 $x\in\mathcal X(k)$, 诱导的稳定子群映射 $G_x\to G_{f(x)}$ 都是单射.

\textbf{代数空间的性质.}
\begin{itemize}
\item (等价关系商的代数性) 若 $R\rightrightarrows U$ 是概形的平展等价关系, 则商层 $U/R$ 是代数空间 (定理~\bookref{4.4.13}).
推论~\bookref{5.5.12}~将此推广到代数空间的光滑等价关系, 推论~\bookref{7.3.6}~将其推广到 fppf 等价关系.
特别地, 代数群在代数空间上的自由作用的商作为代数空间存在.
\item (Zariski 主定理) 将概形的 Zariski 主定理 (\bookref{A.7.3}) 推广, 若 $X\to Y$ 是 Noether 代数空间的拟有限分离态射, 则存在分解 $X\hookrightarrow\widetilde X\to Y$, 其中 $X\hookrightarrow\widetilde X$ 是开浸入, $\widetilde X\to Y$ 有限 (定理~\bookref{5.5.9}).
特别地, $X\to Y$ 拟仿射. Zariski 主定理也适用于 Deligne--Mumford 叠的可表态射, 并在定理~\bookref{7.1.17}~中推广到代数叠的可表态射.
\item (仿射性判据) 由 Serre 仿射性判据 (\bookref{5.5.16}), 代数空间 $X$ 是仿射概形当且仅当 $\Gamma(X,-)$ 在拟凝聚层范畴上正合.
由 Chevalley 仿射性判据 (\bookref{5.5.20}), 若 $X\to Y$ 是 Noether 代数空间的有限满射, 且 $X$ 仿射, 则 $Y$ 也仿射.
\item (代数空间与概形) 若 $X$ 是拟分离代数空间, 则存在稠密开子空间 $U\subseteq X$ 为概形 (定理~\bookref{5.5.1}).
域上局部有限型的拟分离群代数空间是概形 (定理~\bookref{5.5.28}); 特别地, 对角线拟分离的代数叠的域值点的稳定子是群概形.
分离的一维代数空间是概形 (定理~\bookref{5.5.32}).
\end{itemize}

\textbf{Deligne--Mumford 叠的性质.}
\begin{itemize}
\item (群胚商的代数性) 若 $R\rightrightarrows U$ 是概形的平展群胚, 则商叠 $[U/R]$ 是 Deligne--Mumford 叠 (定理~\bookref{4.4.13}).
\item (Deligne--Mumford 叠的局部结构) 若 $\mathcal X$ 是分离的 Deligne--Mumford 叠, $x\in\mathcal X$ 是稳定子为 $G_x$ 的有限型点, 则存在 $x$ 的仿射平展邻域 $[\Spec A/G_x]\to\mathcal X$ (定理~\bookref{5.3.1}).
\item (Keel--Mori 定理) 若 $\mathcal X$ 是 Noether 概形 $S$ 上分离且有限型的 Deligne--Mumford 叠, 则存在粗模空间 $\mathcal X\to X$, 其中 $X$ 是 $S$ 上分离的代数空间 (定理~\bookref{5.4.6}).
\item (Gabber 引理) 若 $\mathcal X$ 是 Noether 概形 $S$ 上分离且有限型的 Deligne--Mumford 叠 (如代数空间), 则存在概形 $Z$ 及有限满射 $Z\to\mathcal X$ (定理~\bookref{5.6.1}).
\end{itemize}

\textbf{代数叠的性质.}
\begin{itemize}
\item (群胚商的代数性) 若 $R\rightrightarrows U$ 是概形的光滑群胚, 则商叠 $[U/R]$ 是代数叠 (定理~\bookref{4.4.13}). 推论~\bookref{7.3.6}~将其推广到代数空间的 fppf 群胚.
\item (剩余 gerbe 与极小呈示) 若 $\mathcal X$ 是 Noether 代数叠, $x\in|\mathcal X|$ 是有限型点, 则\newterm{剩余 gerbe}{residual gerbe} $\mathcal G_x$ 存在, 且 $\mathcal G_x\hookrightarrow\mathcal X$ 是局部闭浸入 (命题~\bookref{4.5.16}).
若进一步要求稳定子 $G_x$ 光滑, 则存在相对维数为 $\dim G_x$ 的光滑呈示 $(U,u)\to(\mathcal X,x)$, 使得 $\mathcal G_x\times_{\mathcal X}U\cong\Spec\kappa(u)$ (定理~\bookref{4.6.1}).
后面将建立任意点 $x\in|\mathcal X|$ 的剩余 gerbe 的存在性 (命题~\bookref{7.4.39}).
\item (无穷小提升判据) 我们将建立无穷小提升判据, 刻画光滑, 平展, 非分歧或对角线非分歧的态射 (定理~\bookref{4.7.1}).
\item (赋值判据) 我们将给出赋值判据, 刻画固有, 普遍闭, 分离或对角线分离的态射 (定理~\bookref{4.8.7}).
\item (代数叠的局部结构) 若 $\mathcal X$ 是代数闭域 $k$ 上具有仿射对角线的有限型代数叠, 则每个具有线性约化稳定子 $G_x$ 的点 $x\in\mathcal X(k)$ 都有仿射平展邻域 $[\Spec(A)/G_x]\to\mathcal X$ (定理~\bookref{7.7.1}).
\item (良模空间的存在性) 设 $\mathcal X$ 是特征 $0$ 的代数闭域 $k$ 上具有仿射对角线的有限型代数叠.
若 $\mathcal X$ 是 $S$-完备且 $\Theta$-完备的, 则存在良模空间 $\mathcal X\to X$, 其中 $X$ 是 $k$ 上分离的有限型代数空间 (定理~\bookref{7.10.1}).
\end{itemize}

\textbf{历史说明}

Deligne--Mumford 叠与代数叠最早在~\cite{DM69,Art74b}~中引入, 两者在那里都称为代数叠, 所用约定与我们略有不同.
具体地说, \cite[定义 $4.6$]{DM69}~假设代数叠具有可由概形表出的平展呈示 $U\to\mathcal X$ (他们在脚注中指出, 当对角线分离且拟紧时, 这给出了``正确的''定义).
另一方面, \cite[定义 $5.1$]{Art74b}~假设代数叠在优良 Dedekind 整环上局部有限型.
我们避免使用的术语 Artin 叠, 有时专指满足 Artin 公理 (\bookref{C.7.1}~或~\bookref{C.7.4}) 的叠, 这些公理最早在~\cite[定理 $5.3$]{Art74b}~中建立.

我们遵循~\cite{LMB00,Ols16,SP}~的约定, 但有以下例外:
\begin{itemize}
\item \cite{LMB00}~对对角线施加分离且拟紧的假设; 我们与~\cite{Ols16,SP}~一样, 不施加任何条件.
\item 我们在 $\Sch_{\text{ét}}$ 上定义代数空间与代数叠, 而~\cite{LMB00}~使用仿射概形的大平展景 $\mathsf{AffSch}_{\text{ét}}$, \cite{SP}~则使用 $\Sch_{\mathrm{fppf}}$.
这些定义给出等价的代数叠概念 (参见~\cite[Tag 076U]{SP}).
\item 对叠的态射, 我们所称的可表 (分别为可由概形表出), 在~\cite[Tags 02ZW 和 02Y7]{SP}~中称为可由代数空间表出 (分别为可表).
\end{itemize}
\flushbottom
